\section{Introduction}
\label{sec:intro}

LLM agents are now deployed with real tool access: they read files, query inboxes and chat logs, and issue outbound HTTP requests on behalf of a user. Most of these tool calls happen in backend pipelines the user never sees. An agent invoked to ``run the daily sync'' silently reads, retrieves, and forwards on the user's behalf, and only a summary returns to the screen. This creates a confused-deputy exposure that the literature has named \emph{indirect prompt injection}: malicious instructions placed in a trusted data or skill channel (an email, a support ticket, a tool's own documentation) are ingested as content and then followed as commands, redirecting the agent's authorized capabilities toward an attacker's ends~\cite{greshake2023}. The injection rides in through exactly the channels the agent must read to do its job.

The defenses the industry reaches for are \emph{behavioral}: a defensive system prompt (an RFC-2119 security policy), an input classifier on the user message, an output filter scanning tool-call parameters against a credential denylist. These are negative, default-allow controls. They enumerate the bad and permit everything else, and they are permanently bypassable for a structural reason rather than a tuning one. A defensive prompt and an input classifier are \emph{in-band}: they live in the same input channel the injection occupies, so an injection that can address the model can address them too, and overrides them. An output filter is a negative denylist, so business data carrying no credential signature sails straight through it. No amount of prompt engineering or classifier retraining removes a weakness that is a property of \emph{where} and \emph{how} the control is applied.

This is the same shape as the long-settled host-security shift from signature-based antivirus (a negative control that blocks known-bad) to application allowlisting (a positive control that permits only known-good). The durable lesson there was that security should be a property of \emph{complete mediation} and \emph{fail-safe defaults}, mediating every access and denying by default, rather than of an ever-growing catalogue of recognized bad~\cite{galloway2020, saltzer1975, anderson1972}. The deeper lineage is older still: object-capability systems make authority an unforgeable token that can only be exercised, never named into existence by an unauthorized party~\cite{dennisvanhorn1966, miller2006}, and the LANGSEC programme makes the recognizer \emph{correct by construction}, so that the accepted language \emph{is} the policy and anything outside it is not merely rejected but ungenerable at the input boundary~\cite{sassaman2013langsec}. We operationalize this lineage for LLM agents. Most directly, this study is a domain transfer of Galloway's controlled evaluation of application allowlisting as a malicious-code protection control~\cite{galloway2020} from host security into AI security: the operating system's executable space becomes the LLM agent's tool-call action space, signature-based detection becomes the behavioral guardrails, and application allowlisting becomes the per-request authorization grammar. The negative controls are the behavioral guardrails above; the positive control is Grammar-Constrained Decoding (GCD) with a \emph{per-request authorization grammar} \(G_s\), an unforgeable capability over the action space~\cite{dennisvanhorn1966, miller2006}, compiled from the agent's authenticated identity and admitting only in-scope tool calls. The attack corpus is held fixed across both. Unlike a structured-output conformance check, every trial is scored by \emph{effect}: a bypass is recorded iff the agent emits a \texttt{fetchURL} call carrying the attack's sentinel marker in \texttt{fetchURL.data} to a non-owned endpoint (\texttt{vendor-compliance.net}, \texttt{okta-verify.net}, or \texttt{*.cybershark-tools.com}), per a deterministic win detector that requires \emph{both} the attacker URL \emph{and} the marker (Section~\ref{sec:method}); data to allowed endpoints does not fire a win. We run this at a scale (millions of attack trials) commensurate with the modern agent-security benchmark landscape, where the Agent Security Bench reports 90{,}000 test cases~\cite{zhang2024asb}, and report every rate as a Clopper--Pearson exact binomial interval~\cite{clopper1934}.

The reason a positive grammar holds where a positive \emph{policy} alone does not is a matter of where enforcement sits. The transformer always computes a full logit distribution, and under injection it may assign an arbitrarily high logit to a forbidden token: the model's ``intent'' is uncontrolled and may be fully adversarial. However, GCD applies the grammar mask \emph{before} sampling, driving every grammar-invalid token's probability to zero, so such a token cannot be generated by any sampler, whether greedy, top-\(k\), top-\(p\), or temperature~\cite{willard2023, dong2024xgrammar}. The model is an untrusted machine, the sampler is the gate, and the grammar is the allowlist. Logits encode intent and are left uncontrolled; the generated tokens are constrained out-of-band, at a point the attacker's tokens never reach.

Beyond security, the same placement yields a second payoff. GCD provides a \emph{validity dividend}: the grammar makes well-formed in-policy tool-call emission a structural property independent of model size. A weak 1.7B under blind GCD emits a valid in-policy call \(\geq\)98.9\% of the time (Section~\ref{sec:res-blast}), while unconstrained Mistral-Small-4-119B produces a valid call only \(\approx\)59\% of the time (Section~\ref{sec:res-mistral}), a model-independent structural property, with correctness about \emph{which} in-scope action left to model intelligence. Here ``valid'' means in-grammar and in-policy (a well-formed authorized call), not semantically correct about which in-scope action to take (the deflection DV); see contribution~4.

We add two layers absent from the host-security empirical tradition. The first is a \emph{soundness theorem}: because the grammar mask quantifies over all logit distributions, the guarantee \(\textrm{emitted} \in L(G_s) \subseteq \mathrm{Authorized}(s)\) holds against a fully injection-compromised model, never assuming the logits are benign. This is what licenses a class-level elimination-by-construction claim that no finite sample could (Section~\ref{sec:formal}). Decode-time enforcement of this kind has direct predecessors. Agent-C proves a temporal-ordering guarantee for agent steps via a decode-time automaton, but does so by an SMT check-and-resample over step ordering rather than a per-principal authorization mask, and its soundness is over ordering, not authority~\cite{agentc2025}. PackMonitor eliminates package hallucinations via the same DFA-plus-logit-mask mechanism, but to a correctness (hallucination) objective rather than authorization~\cite{packmonitor2026}. Our contribution differs in (a) grounding the grammar in the agent's authenticated identity (a per-request capability) and masking the unauthorized token to zero rather than detecting-and-resampling, (b) the constitutive-vs-corrective distinction with the produce-then-catch cost axis below, and (c) operating at the tool-call authorization level across the Tantalus multi-skill agent environment at \({>}2\)~million trials. We claim to be the first to \emph{build and measure} generation-time tool-call authorization at this scale. We do not claim to be the first to pursue least-privilege for tool-calling agents: MiniScope~\cite{miniscope2025}, for instance, reconstructs and enforces per-tool permission hierarchies, but at execution time rather than at generation.

The second added layer is a \emph{cost axis}: token generation is not free. We compare GCD against the strongest deployed positive control, a post-parse allowlist gate (call it the corrective allowlist), the class to which systems such as CaMeL~\cite{camel2025}, Progent~\cite{progent2025}, OAP~\cite{oap2026}, AgentArmor~\cite{agentarmor2025}, and ARM~\cite{arm2026} belong (see Section~\ref{sec:related}). The corrective allowlist enforces the \emph{same} policy over the \emph{same} language \(L(G_s)\), the same programmatic allowlist derived from the grammar's enumerated sets and verified as identical by construction (Section~\ref{sec:method}), but \emph{after} an unbounded generator has already produced the call. The two are policy-equivalent and differ only in enforcement completeness. GCD is \emph{constitutive}: the out-of-scope action is ungenerable and exists as an artifact nowhere (output buffer, logs, telemetry, stream, KV cache). The allowlist is \emph{corrective}: it must generate-then-reject, so the forbidden artifact is produced in full and then caught, re-introducing one layer down the production cost and output-channel artifact (both measured; Section~\ref{sec:results}), as well as, in the general case, a coverage obligation and a TOCTOU window (theoretical properties of the corrective model; Section~\ref{sec:core}). The empirical signature of this difference is the per-trial \emph{emission} count, the number of times the unauthorized artifact actually existed, which is 0 under GCD by construction and \(>0\) under the allowlist.

Why, then, is scale a finding rather than a foregone conclusion? Two reasons. First, bypass is scored by effect through the \emph{unconstrained} surfaces: the free-text fields a grammar cannot enumerate (\texttt{searchFiles.query}, \texttt{respondToUser.message}) remain open, so a 0\% bypass at the structured sink is an empirical claim about what the agent does \emph{with} that latitude, not a restatement of the grammar's membership relation. (The \texttt{c\_sealed} steelman is a constitutive demonstration by a different route: the denylist trajectory is forced and the sink sealed, so its 0\% exfil is substantially by construction for the authorization objective, because the grammar seals the sink deterministically. It is included as an adversarial upper-bound exhibit, not as an independent replication of H2; see Section~\ref{sec:res-steel} and limitations item~3.) Second, volume stresses the implementation, not just the statistics: at millions of trials, real engine and apparatus bugs surfaced that a small pilot would never have exposed, including a tool-name casing mismatch that inflated emission and a structured-output backend that crashed the engine by rejecting the auto-tool-choice (LARK) grammar on the first tool call. Large \(N\) tightens the confidence intervals; it does not, by itself, buy generalization, and we are explicit about that throughout.

We keep the scope honest from the outset. This is a single-subject \emph{feasibility} study on the Tantalus subject system, not a population-level universality claim. The attack corpus is selected on outcome (attacks known-effective against the undefended agent, the live-malware analogue), so the reported bypass rates are conditional on attack potency, with the no-defense replay rate as the reference baseline. The headline security claim is \emph{action integrity over enumerable-sink tools} (structured tool calls and enumerable parameters); free-text confidentiality leakage through channels that must remain free-text is explicitly Paper-2 scope.

\paragraph{Contributions.} This paper makes the following contributions.
\begin{itemize}[leftmargin=1.4em]
  \item \textbf{The constitutive-vs-corrective distinction, and its empirical signature.} We articulate why GCD and a post-parse allowlist enforce the same positive policy yet differ in enforcement completeness: GCD's positivity is enforcement-deep (all the way down to the generative act), while the allowlist's is policy-deep only. We make \emph{emission} the primary structural dependent variable that measures the difference directly (Section~\ref{sec:core}).
  \item \textbf{A soundness theorem quantified over all logit distributions.} We prove \(\textrm{emitted} \in L(G_s) \subseteq \mathrm{Authorized}(s)\) against a fully injection-compromised model, separating the proof's class-level guarantee from the experiment's evidence in a three-tier ladder (Section~\ref{sec:formal}).
  \item \textbf{A large-scale, effect-scored experiment porting the AV-vs-allowlisting design to LLM agents}~\cite{saltzer1975, sassaman2013langsec}. A generator\(\times\)gate factorial design (Section~\ref{sec:method}) over \({\sim}2.46\) million total trials (\(\approx\)\num{2460677}) with Clopper--Pearson intervals: on two capable non-Qwen anchors (Mistral-Small-4-119B and Gemma-4-31B, the second and third of three pretraining lineages; full picture in Section~\ref{sec:results}) the full behavioral stack still leaks (Mistral A4 bypass \(1.82\%\) \([1.43, 2.29]\), Gemma A4 \(33.80\%\) \([31.73, 35.92]\), CI-lower \(>0\)) while GCD holds at \(0.00\%\) (Mistral \([0, 0.09]\), Gemma \([0, 0.18]\)), strictly below every behavioral point estimate; across all GCD attack arms, \(868{,}000\) trials at zero bypass (Section~\ref{sec:results}).
  \item \textbf{Validity decoupled from model capability (the GCD validity dividend).} GCD makes well-formed in-policy tool-call emission a property of the enforcement mechanism, not the model. A 1.7B under GCD emits a valid in-policy call \(\geq\)97.8\% of the time (\texttt{c\_guided}: 97.8\%; \texttt{c}: 98.9\%; C\(+\): 99.9\%); the same grammar raises a capable 119B from 59.2\% (unconstrained) to 99.8\% (GCD). This is a corollary of the soundness theorem (Lemma~\ref{lem:mask-soundness}): any logit distribution's output is constrained to \(L(G_s)\), so well-formedness is guaranteed by construction at the token level. The deployment payoff is separability: capability buys task intelligence, the grammar buys structural validity, so a smaller and cheaper model achieves guaranteed structural validity that a frontier-class model fails to produce \(\approx\)40\% of the time unaided (Section~\ref{sec:results}). The honest fence: validity here means in-grammar and in-policy, so the call never fails \emph{structurally} by construction, but it is not semantically correct about \emph{which} in-scope action (the deflection DV).
  \item \textbf{The cost and reliability axes the empirical tradition lacked.} We quantify the produce-then-catch tax (emission, output tokens, retry attempts, availability stranding) the corrective allowlist pays and GCD avoids, and a least-privilege variant (\texttt{c\_guided}) that turns blind GCD's two-sided reliability story into a one-sided utility win (Sections~\ref{sec:results} and~\ref{sec:discussion}). However, the registered 2~pp non-inferiority test for blind C on legitimate tasks, although run at power (Section~\ref{sec:res-ladder}), does not establish non-inferiority (CI half-width \({\sim}5.3\)~pp \(\gg\) 2~pp margin); the utility win is established only for the least-privilege variant.
\end{itemize}

\paragraph{Roadmap.} Section~\ref{sec:threat} fixes the threat model and the effect-scored bypass definition; Section~\ref{sec:core} develops the constitutive-vs-corrective core; Section~\ref{sec:formal} states and proves the soundness theorem and the three-tier ladder; Section~\ref{sec:method} describes the subject system, conditions, and factorial design; Section~\ref{sec:results} reports the five datasets (Sections~\ref{sec:res-blast}, \ref{sec:res-mistral}, \ref{sec:res-ladder}, \ref{sec:res-gemma}, \ref{sec:res-steel}) and the aggregate hard-statistics ledger (Section~\ref{sec:res-ledger}). Sections~\ref{sec:discussion} and~\ref{sec:limits} discuss implications and limitations honestly, with related work in Section~\ref{sec:related} and conclusions in Section~\ref{sec:conclusion}.
